\section{Jina v5 paired response
similarity}\label{record-docs-jina-similarity-md-jina-v5-paired-response-similarity}

This adds the same embedding-based metric to ordinary paraphrasing,
BIRA, direct probability resampling and confidence smoothing for both
the TextSeal 2 and DistSeal-inspired instruct-Qwen experiments. It does
not alter their attacks or select outputs using owner detections.

The pinned model is the official
\href{https://huggingface.co/jinaai/jina-embeddings-v5-text-small-text-matching}{jinaai/jina-embeddings-v5-text-small-text-matching}
checkpoint \texttt{\seqsplit{69890c317b81d4e1f412e66ebe68b6ebf9534d24}},
which merges Jina v5's text-matching adapter into its backbone. The
released Sentence Transformers configuration specifies the
\texttt{\seqsplit{Document:}} prefix, last non-padding token pooling,
full 1024-dimensional embeddings and L2 normalization. The evaluator
uses the same operations with native Transformers and SDPA; it does not
require remote code execution or an external API. No evaluation text
leaves the research VM.

Original and candidate responses are aligned by prompt ID, with missing
or duplicate IDs rejected. The original response is the unmodified
marked TextSeal 2 output or the unmodified distilled student's response.
Each method has 256 paired test questions. All inputs and code are
frozen and hashed. Jina inference uses no truncation; inputs beyond its
32,768-token limit cause an explicit error. Vectors are cached by text
hash in restartable batches.

For paired embeddings a and b, the score is their cosine similarity:
\texttt{\seqsplit{dot(a, b) / (norm(a) * norm(b))}}. Scores are reported
directly, not converted into a percentage or a preservation pass rate.
Each condition reports mean, median, standard deviation, P05/P25/P75/P95
and a 95\% bootstrap interval for the mean (2,000 resamples over
question pairs). Paired differences against ordinary paraphrasing use
the same question-level sampling unit.

Controls include original-versus-itself, original-versus-clean-model
answers, original-versus-human answers, and a deterministic permutation
pairing different questions. The DistSeal case also includes the learned
teacher's response. Empty pairs are flagged and excluded from the
primary score summary.

For resampling and smoothing, this measures similarity between two
answers to the same question. It is not a requirement that independent
valid answers say exactly the same things. High embedding similarity can
also coexist with wrong numbers, negation or other factual changes; NLI
and factual checks remain useful. No universal cosine threshold is
inferred from these data.

\subsection{Runtime and
commands}\label{record-docs-jina-similarity-md-runtime-and-commands}

The existing Torch 2.9 paper runtime is reused without modification. A
separate \texttt{\seqsplit{jina\_v5\_overlay}} supplies Transformers
5.17 and compatible tokenizers, avoiding an upgrade of either running
experiment environment. Pins are in
\texttt{\seqsplit{requirements\_jina\_similarity.in}}. Downloaded
weights and their revision are recorded in
\texttt{\seqsplit{results/research\_workflow/watermarking/resources/jina\_v5\_model.json}}.

Download the pinned model snapshot with the Hugging Face client and
record its local path before preparation:

\begin{Verbatim}
import json
from pathlib import Path
from huggingface_hub import snapshot_download
resources = Path("results/research_workflow/watermarking/resources")
model = "jinaai/jina-embeddings-v5-text-small-text-matching"
revision = "69890c317b81d4e1f412e66ebe68b6ebf9534d24"
path = snapshot_download(model, revision=revision, cache_dir=str(resources/"hf"),
    allow_patterns=["*.json", "*.txt", "*.safetensors", "*.md", "*.jinja", "1_Pooling/*"])
(resources/"jina_v5_model.json").write_text(json.dumps(
    {"name": model, "revision": revision, "path": str(Path(path).resolve())}))
\end{Verbatim}

\begin{Verbatim}
uv pip install --python results/research_workflow/watermarking/resources/bira_runtime_py310/bin/python \
  --target results/research_workflow/watermarking/resources/jina_v5_overlay \
  --no-deps -r requirements_jina_similarity.in
PYTHONPATH=. uv run --no-sync python scripts/jina_watermark_similarity.py prepare TEXT_RUN/manifest.json
PYTHONPATH=. uv run --no-sync python scripts/jina_watermark_similarity.py submit JINA_RUN/manifest.json
PYTHONPATH=. uv run --no-sync python -m pytest ai_content_detector/tests/test_semantic_similarity.py -q
\end{Verbatim}

The TextSeal 2 job can run immediately once its producer is complete.
The DistSeal job waits for the existing teacher/student/attack job to
finish, then freezes its final attack artifacts. Each uses one H100 on
two-hour dev QoS, saves every embedding batch and requeues on the
advance signal. Project storage holds the job caches, which are cleaned
on exit.

Five CPU checks passed in this checkout: ID alignment, mixed-padding
last-token pooling, cosine normalization, pair-level bootstrap behavior
and empty-pair flags. GPU evaluation outcomes are recorded per case in
\texttt{\seqsplit{receipt.json}}, \texttt{\seqsplit{summary.json}} and
\texttt{\seqsplit{report.md}}; job submission alone is not a completed
comparison.

Both GPU evaluations completed successfully in run
\texttt{\seqsplit{20261104T081330240270Z}} (jobs
\texttt{\seqsplit{1938571}} and \texttt{\seqsplit{1938573}}). The
real-weight CPU smoke check also passed. See the
\href{records/docs/jina_similarity_results.tex}{saved comparison table}
for the results.
